\subsection{Partial and Total Derivatives}

% **Role:** calculus as local linearization plus global geometric integration.
%
% Linear approximation, Jacobians, partial/directional derivatives, chain rule.

Several-variable differentiation treats derivatives as best linear approximations, with operator norms controlling the remainder.\\
Determinants and invertibility encode the resulting local geometry.

\begin{definition} \label{def:total-derivative-rn-rm} \hlt{Total Derivative}\\
Let $E\subseteq\R^n$ be open, $\f:E\to\R^m$, $\x\in E$. Write $\LL(\R^n,\R^m)$ for space of linear maps from $\R^n$ to $\R^m$ (Definition \ref{def:linear-map}).\\
The map $\f$ is \hlt{differentiable at $\x$} iff $(\exists\,A\in\LL(\R^n,\R^m))\,\lim\limits_{\h\to\ze}\frac{\norm{\f(\x+\h)-\f(\x)-A\h}_2}{\norm{\h}_2}=0$.\\
Any such $A$ is a \hlt{total derivative} of $\f$ at $\x$; when it exists, it is denoted $D\f(\x)$ or $\f'(\x)$.\\
Equivalently, with $\mathbf r(\h)\defeq\f(\x+\h)-\f(\x)-D\f(\x)\h$, one has $\norm{\mathbf r(\h)}_2=o(\norm{\h}_2)$.\\
The map $\f$ is \hlt{differentiable on $E$} iff $(\forall\,\x\in E)\,D\f(\x)$ exists.
\end{definition}

\begin{perspective} \label{per:total-derivative-linear-map-first} \hlt{Derivative First, Matrix Second}\\
The total derivative is a linear map by Definition \ref{def:total-derivative-rn-rm}; its coordinate matrix is secondary.\\
The intrinsic control is $(\forall\,A\in\LL(\R^n,\R^m))(\forall\,\h\in\R^n)\,\norm{A\h}_2\leq\normop{A}\norm{\h}_2$ by Definition \ref{def:operator-norm-rn-rm}.
\end{perspective}

\noindent Figure \ref{fig:total-derivative-linear-approximation} separates the input increment $\h$ from its image $D\f(\x)\h$. The dashed vector is the smaller remainder $\mathbf r(\h)$.
\begin{illustrationfigure}[H]
  \centering
  \begin{tikzpicture}[
  scale=0.88,
  >=Stealth,
  every node/.style={font=\scriptsize},
  axis/.style={->, gray!70, line width=0.45pt},
  vector/.style={->, black, line width=0.75pt},
  linear/.style={->, green!45!black, line width=0.85pt},
  error/.style={->, black!65, dashed, dash pattern=on 3pt off 2pt, line width=0.55pt},
  point/.style={circle, fill=black, inner sep=1.15pt},
  approxpoint/.style={circle, fill=green!45!black, inner sep=1.15pt}
]
  \begin{scope}
    \draw[axis] (-0.25,0) -- (3.0,0) node[right] {$x$};
    \draw[axis] (0,-0.2) -- (0,2.1) node[above] {$y$};

    \coordinate (X) at (0.85,0.65);
    \coordinate (XH) at (2.25,1.25);

    \node[point] at (X) {};
    \node[point] at (XH) {};
    \draw[vector] (X) -- node[above left=-1pt] {$\h$} (XH);

    \node[below left=1pt] at (X) {$\x$};
    \node[right] at (XH) {$\x+\h$};
    \node[above] at (1.35,2.25) {domain};
  \end{scope}

  \begin{scope}[xshift=4.2cm]
    \draw[axis] (-0.25,0) -- (3.45,0) node[right] {$x$};
    \draw[axis] (0,-0.2) -- (0,2.35) node[above] {$y$};

    \coordinate (FX) at (0.65,0.72);
    \coordinate (LIN) at (2.18,1.45);
    \coordinate (ACT) at (2.98,1.95);

    \node[point] at (FX) {};
    \node[approxpoint] at (LIN) {};
    \node[point] at (ACT) {};

    \draw[linear] (FX) -- node[below right=-1pt] {$D\f(\x)\h$} (LIN);
    \draw[error] (LIN) -- node[below right=-1pt] {$\mathbf r(\h)$} (ACT);

    \node[left] at (FX) {$\f(\x)$};
    \node[above left=2pt] at (LIN) {$\f(\x)+D\f(\x)\h$};
    \node[above right=1pt] at (ACT) {$\f(\x+\h)$};
    \node[above] at (1.6,2.25) {codomain};
  \end{scope}
\end{tikzpicture}

  \caption{The solid vector $D\f(\x)\h$ reaches the linear approximation; the dashed $\mathbf r(\h)=o(\norm{\h}_2)$ completes the actual increment.}
  \label{fig:total-derivative-linear-approximation}
\end{illustrationfigure}

\begin{theorem} \label{thm:total-derivative-unique-and-continuous} \hlt{Basic Consequences of Total Differentiability}\\
Let $E\subseteq\R^n$ be open, let $\f:E\to\R^m$, and let $\x\in E$. If $\f$ is differentiable at $\x$, then:
\begin{enumerate}[label=(\roman*)]
\setlength{\itemsep}{0pt}
\item the total derivative $D\f(\x)$ is unique;
\item $\f$ is continuous at $\x$.
\end{enumerate}
\begin{proof}
\begin{enumerate}[label=(\roman*)]
\setlength{\itemsep}{0pt}
\item Let $A_1,A_2$ be total derivatives at $\x$, $B\defeq A_1-A_2$. If $\mathbf r_j(\h)\defeq\f(\x+\h)-\f(\x)-A_j\h$, then $B\h=\mathbf r_2(\h)-\mathbf r_1(\h)$.\\
$\norm{B\h}_2=o(\norm{\h}_2)$ by Definition \ref{def:total-derivative-rn-rm},  triangle inequality. For $\mathbf u\neq\ze$ and $t\neq0$, $\frac{\norm{B(t\mathbf u)}_2}{\norm{t\mathbf u}_2}=\frac{\norm{B\mathbf u}_2}{\norm{\mathbf u}_2}$ by linearity,  homogeneity.\\
For $\mathbf u\neq\ze$, $\norm{B\mathbf u}_2=0$ by this equality and its limit as $t\to0$; $B\ze=\ze$ by linearity. Hence $B=0$.
\item Set $A\defeq D\f(\x)$, $\mathbf r(\h)\defeq\f(\x+\h)-\f(\x)-A\h$. Then $\norm{\f(\x+\h)-\f(\x)}_2\leq\normop{A}\norm{\h}_2+\norm{\mathbf r(\h)}_2\to0$ by Definitions \ref{def:operator-norm-rn-rm}, \ref{def:total-derivative-rn-rm}.\\
Thus $\f(\x+\h)\to\f(\x)$ as $\h\to\ze$, so $\f$ is continuous at $\x$.
\end{enumerate}
\hlt{Structural Commentary}: Scaling eliminates a second linear part; the operator bound transfers the remainder estimate to continuity.\\
\hlt{Key Insight}: A nonzero linear error cannot be $o(\norm{\h}_2)$.
\end{proof}
\end{theorem}

\begin{example} \label{ex:derivative-of-linear-map} \hlt{Derivative of a Linear Map}\\
Let $A\in\LL(\R^n,\R^m)$, $\f(\x)\defeq A\x$. Then $(\forall\x,\h\in\R^n)\,\f(\x+\h)-\f(\x)-A\h=\ze$, so $(\forall\x\in\R^n)\,D\f(\x)=A$ by Definition \ref{def:total-derivative-rn-rm}.
\end{example}

\begin{theorem} \label{thm:chain-rule-several-variables} \hlt{Chain Rule}\\
Let $E\subseteq\R^n$ and $G\subseteq\R^m$ be open, $\x_0\in E$. Let $\f:E\to\R^m$, $\g:G\to\R^k$, with $\f(E)\subseteq G$.\\
If $\f$ is differentiable at $\x_0$ and $\g$ is differentiable at $\f(\x_0)$, then:
\begin{enumerate}[label=(\roman*)]
\setlength{\itemsep}{0pt}
\item $\g\circ\f$ is differentiable at $\x_0$;
\item $D(\g\circ\f)(\x_0)=D\g(\f(\x_0))D\f(\x_0)$.
\end{enumerate}
\begin{proof}
Set $A\defeq D\f(\x_0)$, $B\defeq D\g(\f(\x_0))$. Then $BA\in\LL(\R^n,\R^k)$ by Proposition \ref{prop:linear-map-product-algebra}(i).\\
Set $\mathbf u(\h)\defeq\f(\x_0+\h)-\f(\x_0)-A\h$, $\mathbf q(\h)\defeq\f(\x_0+\h)-\f(\x_0)$, $\mathbf v(\mathbf q)\defeq\g(\f(\x_0)+\mathbf q)-\g(\f(\x_0))-B\mathbf q$.
\begin{enumerate}[label=(\roman*)]
\setlength{\itemsep}{0pt}
\item Relation $\norm{\mathbf u(\h)}_2=o(\norm{\h}_2)$ holds by Definition \ref{def:total-derivative-rn-rm}. Relation $\norm{\mathbf v(\mathbf q)}_2=o(\norm{\mathbf q}_2)$ as $\mathbf q\to\ze$ holds by same definition.\\
Eventually $\norm{\mathbf u(\h)}_2\leq\norm{\h}_2$ as $\h\to\ze$ by the first estimate.\\
Thus $\norm{\mathbf q(\h)}_2\leq(\normop{A}+1)\norm{\h}_2$ eventually by $\mathbf q(\h)=A\h+\mathbf u(\h)$, Definition \ref{def:operator-norm-rn-rm}; $\norm{\mathbf v(\mathbf q(\h))}_2=o(\norm{\h}_2)$.\\
The identity $(\g\circ\f)(\x_0+\h)-(\g\circ\f)(\x_0)-BA\h=B\mathbf u(\h)+\mathbf v(\mathbf q(\h))$ holds by the definitions of $\mathbf u$, $\mathbf q$, and $\mathbf v$.\\
The right side is $o(\norm{\h}_2)$ by the preceding estimates and the operator-norm bound, so $\g\circ\f$ is differentiable at $\x_0$.
\item The derivative candidate in (i) is $BA$. Thus $D(\g\circ\f)(\x_0)=BA=D\g(\f(\x_0))D\f(\x_0)$ by Theorem \ref{thm:total-derivative-unique-and-continuous}(i).
\end{enumerate}
\hlt{Structural Commentary}: The inner increment is $O(\norm{\h}_2)$, allowing the outer remainder to compose with it.\\
\hlt{Key Insight}: Both composed error terms remain $o(\norm{\h}_2)$.
\end{proof}
\end{theorem}

\begin{theorem} \label{thm:several-variable-componentwise-differentiability} \hlt{Componentwise Differentiability}\\
Let $E\subseteq\R^n$ be open, $\f=(f_1,\ldots,f_m):E\to\R^m$, $\x\in E$. For $i\in J_m$, let $\pi_i\in\LL(\R^m,\R)$ be the $i$th coordinate projection.
\begin{enumerate}[label=(\roman*)]
\setlength{\itemsep}{0pt}
\item $\f$ is differentiable at $\x$ iff $(\forall\,i\in J_m)\,f_i$ is differentiable at $\x$.
\item If $\f$ is differentiable at $\x$, then $(\forall\,i\in J_m)\,D f_i(\x)=\pi_i\circ D\f(\x)$.
\end{enumerate}
\begin{proof}
\begin{enumerate}[label=(\roman*)]
\setlength{\itemsep}{0pt}
\item[$(i)$] ($\Rightarrow$) Let $\mathbf r(\h)\defeq\f(\x+\h)-\f(\x)-D\f(\x)\h$. For $i\in J_m$, set $r_i(\h)\defeq\pi_i(\mathbf r(\h))$. Map $\pi_i\circ D\f(\x)$ is linear by Proposition \ref{prop:linear-map-product-algebra}(i).\\
Then $f_i(\x+\h)-f_i(\x)-(\pi_i\circ D\f(\x))\h=r_i(\h)$ and $\abs{r_i(\h)}\leq\norm{\mathbf r(\h)}_2=o(\norm{\h}_2)$.\\
Thus all component functions are differentiable at $\x$ by Definition \ref{def:total-derivative-rn-rm}.
\item[$(i)$] ($\Leftarrow$) Suppose $(\forall\,i\in J_m)\,f_i$ is differentiable at $\x$,   $A_i\defeq D f_i(\x)$. Define $A:\R^n\to\R^m$ by $A\h\defeq(A_1\h,\ldots,A_m\h)$.\\
$A\in\LL(\R^n,\R^m)$ by coordinatewise linearity. For $i\in J_m$, set $r_i(\h)\defeq f_i(\x+\h)-f_i(\x)-A_i\h$. Then $(\forall\,i\in J_m)\,\abs{r_i(\h)}=o(\norm{\h}_2)$.\\
For $\mathbf r(\h)\defeq(r_1(\h),\ldots,r_m(\h))$, $\norm{\mathbf r(\h)}_2\leq\sum\limits_{i=1}^{m}\abs{r_i(\h)}=o(\norm{\h}_2)$ by finiteness of $J_m$.\\
Hence $\f$ is differentiable at $\x$ with derivative $A$ by Definition \ref{def:total-derivative-rn-rm}.
\item[$(ii)$] For $i\in J_m$, derivative candidate in $(i)$, ($\Rightarrow$), is $\pi_i\circ D\f(\x)$; $(\forall\,i\in J_m)\,D f_i(\x)=\pi_i\circ D\f(\x)$ by Theorem \ref{thm:total-derivative-unique-and-continuous}(i).
\end{enumerate}
\hlt{Structural Commentary}: Coordinate projections prove necessity; finite assembly of scalar remainders proves sufficiency.\\
\hlt{Key Insight}: Finitely many scalar $o(\norm{\h}_2)$ terms form one vector remainder.
\end{proof}
\end{theorem}

\begin{definition} \label{def:partials-gradient-directional-derivatives} \hlt{Partial, Gradient, and Directional Derivatives}\\
Let $E\subseteq\R^n$ be open, $\f=(f_1,\ldots,f_m):E\to\R^m$, $\x\in E$. Let $\mathbf e_1,\ldots,\mathbf e_n$ be standard basis of $\R^n$.
\begin{enumerate}[label=(\roman*)]
\setlength{\itemsep}{0pt}
\item For $(i,j)\in J_m\times J_n$, set $\phi_{ij}(t)\defeq f_i(\x+t\mathbf e_j)$ if $\x+t\mathbf e_j\in E$. When limit exists, \hlt{partial derivative} is $(D_jf_i)(\x)\defeq\lim\limits_{t\to0}\frac{\phi_{ij}(t)-\phi_{ij}(0)}{t}$.
\item For $\mathbf u\in\R^n$, set $\gamma_{\mathbf u}(t)\defeq\f(\x+t\mathbf u)$ if $\x+t\mathbf u\in E$. When limit exists, \hlt{directional derivative} is $(D_{\mathbf u}\f)(\x)\defeq\lim\limits_{t\to0}\frac{\gamma_{\mathbf u}(t)-\gamma_{\mathbf u}(0)}{t}$.\\
When only direction matters, take $\norm{\mathbf{u}}_2=1$.
\item If $m=1$, write $f=f_1$; if all $(D_jf)(\x)$ exist, the \hlt{gradient} of $f$ at $\x$ is $(\nabla f)(\x)\defeq\sum\limits_{j=1}^{n}(D_jf)(\x)\mathbf{e}_j$.
\end{enumerate}
\end{definition}

\begin{perspective} \label{per:partials-are-coordinate-tests} \hlt{Partials Are Coordinate Tests}\\
Partial and directional derivatives in Definition \ref{def:partials-gradient-directional-derivatives} test selected lines.\\
The total derivative in Definition \ref{def:total-derivative-rn-rm} controls every small increment by one linear map.\\
The matrix theorem below extracts the line tests; the next subsection proves a continuity-based converse.
\end{perspective}

\begin{theorem} \label{thm:total-derivative-matrix-partials} \hlt{Matrix of the Total Derivative}\\
Let $E\subseteq\R^n$ be open, $\f=(f_1,\ldots,f_m):E\to\R^m$, $\x\in E$. If $\f$ is differentiable at $\x$, then:
\begin{enumerate}[label=(\roman*)]
\setlength{\itemsep}{0pt}
\item all $(D_jf_i)(\x)$ exist, and $(\forall\,(i,j)\in J_m\times J_n)\,[D\f(\x)]_{ij}=(D_jf_i)(\x)$;
\item if $m=1$ and $f=f_1$, then $(\forall\,\mathbf u\in\R^n)\,(D_{\mathbf u}f)(\x)=(\nabla f)(\x)\cdot\mathbf u$.
\end{enumerate}
\begin{proof}
Set $\mathbf r(\h)\defeq\f(\x+\h)-\f(\x)-D\f(\x)\h$; then $\norm{\mathbf r(\h)}_2=o(\norm{\h}_2)$ by Definition \ref{def:total-derivative-rn-rm}.
\begin{enumerate}[label=(\roman*)]
\setlength{\itemsep}{0pt}
\item For $j\in J_n$, set $\h=t\mathbf e_j$. Then $\frac{\f(\x+t\mathbf e_j)-\f(\x)}{t}=D\f(\x)\mathbf e_j+\frac{\mathbf r(t\mathbf e_j)}{t}\to D\f(\x)\mathbf e_j$ as $t\to0$ by Definition \ref{def:total-derivative-rn-rm}.\\
Thus all component limits exist, and their values form column $j$ of $[D\f(\x)]$ by Proposition \ref{prop:matrix-representation-action}.
\item Suppose $m=1$, $\mathbf u=(u_1,\ldots,u_n)\in\R^n$. Identity $(D_{\mathbf u}f)(\x)=Df(\x)\mathbf u$ holds by Definition \ref{def:total-derivative-rn-rm} with $\h=t\mathbf u$.\\
By part (i) and linearity, $Df(\x)\mathbf u=\sum\limits_{j=1}^{n}u_jDf(\x)\mathbf e_j=\sum\limits_{j=1}^{n}u_j(D_jf)(\x)=(\nabla f)(\x)\cdot\mathbf u$.
\end{enumerate}
\hlt{Structural Commentary}: Line restrictions identify $D\f(\x)\mathbf e_j$; matrix coordinates and linearity recover both claims.\\
\hlt{Key Insight}: The standard basis converts derivative action into partial-derivative columns.
\end{proof}
\end{theorem}

\begin{definition} \label{def:jacobian-determinant} \hlt{Jacobian}\\
Let $E\subseteq\R^n$ be open, $\f=(y_1,\ldots,y_n):E\to\R^n$, $\f$ is differentiable at $\x\in E$.\\
The \hlt{Jacobian determinant} of $\f$ at $\x$ is $J_{\f}(\x)\defeq\det D\f(\x)=\frac{\partial(y_1,\ldots,y_n)}{\partial(x_1,\ldots,x_n)}(\x)$.
\end{definition}

\begin{corollary} \label{cor:jacobian-nonzero-iff-derivative-invertible} \hlt{Nonvanishing Jacobian}\\
Let $E\subseteq\R^n$ be open, $\f:E\to\R^n$ be differentiable at $\x\in E$, $J_{\f}(\x)$ be its Jacobian determinant. Then $D\f(\x)$ is invertible iff $J_{\f}(\x)\neq0$.
\begin{proof}
\begin{enumerate}[label=\roman*.]
\setlength{\itemsep}{0pt}
\item[$(\Rightarrow)$] If $D\f(\x)$ is invertible, then $\det D\f(\x)\neq0$ by Theorem \ref{thm:determinant-basic-laws}(vii). Hence $J_{\f}(\x)\neq0$ by Definition \ref{def:jacobian-determinant}.
\item[$(\Leftarrow)$] If $J_{\f}(\x)\neq0$, then $\det D\f(\x)\neq0$ by Definition \ref{def:jacobian-determinant}. Hence $D\f(\x)$ is invertible by Theorem \ref{thm:determinant-basic-laws}(vii).
\end{enumerate}
\hlt{Structural Commentary}: The Jacobian definition transfers the determinant criterion to the derivative matrix.\\
\hlt{Key Insight}: Nonzero determinant is exactly linear invertibility.
\end{proof}
\end{corollary}

\noindent For the map in Definition \ref{def:jacobian-determinant}, $\abs{J_{\f}(\x)}$ is the volume scale of the linear approximation $D\f(\x)$.
